\section{Related work}

\paragraph{Welfare and optimal-value inference.}
The estimand is the population welfare attained by choosing the better treatment in each stratum. It corresponds to the unrestricted binary first-best benchmark in \citet[Section~2, equations~(3)--(7)]{Cerulli2026OptimalPolicy}. Statistical treatment choice and policy learning evaluate data-dependent treatment rules through their welfare and regret \citep{Manski2004,Kitagawa2018,Athey2021}. Here the decision criterion concerns precision in estimating the scalar population optimum. For the \leanref{S-1}{randomized binary trial with public categorical strata} studied here, the causal model combines uniform stratum probabilities, fair assignment, consistency, and interior arm means, with independent participant records \cref{obj:def:causal-model,obj:ass:iid-people}. The welfare decomposition in \cref{obj:prop:signed-causal-bridge} connects this optimum to the baseline outcome mean and the average absolute treatment contrast. Exact equality of the arm means is therefore part of the parameter space over which precision is evaluated.

Optimal-value inference has developed several approaches to the resulting nonsmoothness. \citet{LuedtkeVanDerLaan2016OptimalValue} characterize pathwise differentiability and construct asymptotic inference for a possibly non-unique optimal treatment strategy, including inference at exceptional laws under their stated conditions. Softmax smoothing yields asymptotic inference under controlled behavior of treatment gaps near zero, suitable smoothing sequences, and nuisance-estimation conditions \citep[Theorem~3.3]{WhitehouseChenAusternSyrgkanis2026Softmax}. For optimal welfare gains integrated against a known target density, \citet[Theorem~1, Assumptions~1, 2(b)(c), and~3]{Chen2025} obtain asymptotic normality under smooth regression functions, overlap, a bounded target-to-training density ratio, a regular zero level set, and the specified uniform first-stage convergence rate. \citet[Sections~3.2 and~4]{Feng2026} develop differentiability results for sorting operators and asymptotic inference for allocation values under their level-set regularity and estimation conditions. Adaptive smoothing and efficiency results in \citet[Theorems~2.2 and~3.1, Assumptions~1--8]{Xu2026} accommodate nonregularity under the stated overlap, moment, margin, nuisance-rate, bandwidth, and homoscedasticity conditions. These contributions establish condition-specific asymptotic guarantees, with several explicitly accommodating treatment ties. Our comparison concerns finite-resource precision uniformly over the declared categorical causal class, including exact ties, and connected confidence intervals with coverage at least \(0.90\) at every law \cref{obj:def:risk,obj:def:honest-length}.

\paragraph{Nonsmooth estimation and private moments.}
Polynomial approximation and moment matching provide the closest statistical precedents. \citet[Theorem~1 and Sections~3--5]{CaiLow2011Nonsmooth} combine a composite-testing lower-bound inequality with Gaussian absolute-mean estimation using polynomial approximation, unbiased Hermite-polynomial estimates, and a hybrid construction. For independent unit-variance Gaussian observations with a fixed coordinate bound, their minimax mean-squared error is asymptotic to \(\beta_*^2M^2[\log\log d/\log d]^2\). The same squared log-log over log form appears here when effective private resource \(t=n\varepsilon^2\) is proportional to \(d^2\), although private transcripts replace Gaussian coordinates and the present result is uniform over the full resource domain. In the discrete sampling model, \citet[Theorems~2--4]{JiaoHanWeissman2018L1} obtain worst-reference mean-squared error of order \(S/(n\log n)\) for raw-sample \(L_1\)-distance estimation under their large-alphabet comparability conditions. That setting observes categorical draws directly; the present paired problem estimates a population distance from locally private original-record messages.

The present construction implements this approximation strategy through participant-level private messages. Products over distinct participants furnish unbiased estimates of contrast powers \cref{obj:def:private-moments}. Their use in polynomial estimation requires bounds that account for dependence among coordinates released in the same messages \cref{obj:lem:private-moment-and-dependence,obj:lem:finite-private-polynomial-calibration}. On the lower-bound side, moment-matched priors are compared through the distribution of the entire private transcript. The sequential contraction result controls their distinguishability when each participant's release can depend on preceding messages \cref{obj:lem:adaptive-moment-contraction}. Thus the approximation comparison is attached to two features of the observation model: attainable moments arise from distinct-person private products, and the converse applies after history-dependent privatization.

\paragraph{Confidentiality and sequential interaction.}
Under \leanref{S-2}{pure local differential privacy}, the target functional and the permitted interaction jointly determine statistical precision. For positive power sums with exponent \(\gamma>1\), \citet[Theorems~2.4--2.7]{ButuceaIssartel2021Functionals} give sequential mean-squared error bounded at order \(t^{-((\gamma-1)\wedge1)}\), with matching \(t^{-1}\) converses for \(\gamma\ge2\). For integrated squared densities over their bounded Besov classes, \citet[Theorems~3.2--3.3 and~4.1--4.2]{ButuceaRohdeSteinberger2023Quadratic} obtain noninteractive power \(t^{-8s/(4s+3)}\) and sequential power \(t^{-4s/(2s+1)}\), up to their stated logarithmic factors; their parametric upper regimes begin at \(s>3/4\) and \(s>1/2\), respectively. The functional here averages absolute signed contrasts, whose kink occurs at equality of the treatment means. Conditional on the cited, unformalized Bernstein absolute-approximation limit, the matching noninteractive construction and adaptive sequential converse establish the same minimax rates in both protocol classes for the specified causal and paired models \cref{obj:thm:uniform-private-value-frontiers,obj:thm:paired-uniform-tv-frontier}. The extra work is to control dependence among private coordinate moments and to contract higher-order, moment-matched transcript mixtures uniformly over adaptive histories and resources.

Average treatment effects provide a complementary causal comparison. \citet[Lemma~4 and Theorem~7]{OhnishiAwan2025Randomized} give an order-\(t^{-1}\) known-propensity estimator based on a bounded inverse-probability score, a matching private minimax comparison, and asymptotic intervals. Their average treatment effect is a signed mean contrast; optimal welfare incorporates the average magnitude of the stratum-specific contrasts. The confidentiality condition here protects the complete observed record, including stratum membership, assignment, and outcome \cref{obj:ass:full-record-privacy}.

The interval criterion also connects to honest-confidence theory. \citet{Low1997} studies expected-length limits under uniform coverage, \citet{CaiLow2004AdaptationCI} develops modulus-based expected-length theory for linear functionals, and \citet{ArmstrongKolesar2020HonestCI} constructs bias-aware nonparametric intervals. Here honesty is finite-sample and global over the declared nonlinear causal or paired model, performance is worst-law expected connected length, and the upper interval is obtained from the uniform squared-error construction. The lower bound is interval-specific: connectedness forces a single reported set to span separated target neighborhoods when private transcript mixtures are difficult to distinguish.

Public randomness also matters for private identity testing: \citet{AcharyaCanonneFreitagTyagi2019TestWithoutTrust} and \citet{Acharya2021} establish its benefits for distinguishing a reference distribution from separated alternatives. Our paired-model criterion estimates the numerical distance from the public uniform reference and assesses globally honest connected intervals for that distance \cref{obj:def:tv-private-decision-values}. A further comparison concerns empirical distribution sketches. \citet[Theorems~4.6, 4.10, and~4.13]{CormodeTing2025FederatedShift} provide sketch accuracy guarantees and approximate item-level local privacy for empirical distribution shifts in the federated setting. The present population criteria use pure privacy for each complete participant record and allow public randomness within the admissible sequential protocols.

Conditional on the cited, unformalized Bernstein absolute-approximation limit, the contribution is the joint finite-resource characterization, up to universal constants, of minimax squared error and minimax worst-law expected length of globally honest connected intervals for the specified causal and paired signed models. Noninteractive procedures attain these bounds, while the converses cover every admissible one-message sequentially interactive protocol \cref{obj:thm:uniform-private-value-frontiers,obj:thm:paired-uniform-tv-frontier}. Under the same premise, embedding the paired family in the full probability simplex also yields the corresponding full-simplex estimation and honest-length lower bounds \cref{obj:thm:full-simplex-tv-converse}.
